\documentclass[11pt]{article}
\usepackage[margin=0.8in]{geometry}
\usepackage{booktabs}
\usepackage{float}
\usepackage{graphicx}
\usepackage{adjustbox}
\usepackage{caption}
\usepackage{array}
\usepackage{hyperref}

\title{Distributional Empirical Discipline: Data and Current Model}
\author{}
\date{Built May 8, 2026}

\begin{document}
\maketitle

\noindent This note is a discipline file, not a validation claim. For each empirical object that is currently available, it shows the data object and the closest object in the active Python model. The model is the current direct-geometry benchmark, job 37, evaluation 1425, with the renewal-valve stationary closure. The benchmark loss in the cluster record is 8.128 and the resolved equilibrium error is 0.0030.

\medskip
\noindent The organizing economic references are standard: Becker and Becker--Lewis for the fertility quantity-quality tradeoff; Sommer and Doepke--Kindermann style lifecycle fertility models for income and lifecycle discipline; Rosen--Roback, Henderson--Ioannides, Sommer--Sullivan--Verbrugge, and Gyourko--Mayer--Sinai for spatial sorting, tenure, and housing-price discipline. The point of this file is to make clear which of those cross-sectional facts the current model can speak to now.

\medskip
\noindent Income needs special care. The model has a deterministic lifecycle income profile, so income moves with age. It does not have persistent idiosyncratic income risk or a within-age income distribution. Therefore the model income-bin plot below uses the stationary mass generated by deterministic age/location income; it is a lifecycle-rank diagnostic, not a clean analog of PSID young-adult income quintiles.


\begin{table}[H]
\centering
\caption{Benchmark calibration moments, with targets next to model moments.}
\label{tab:calibration}
\scriptsize
\begin{adjustbox}{max width=\textwidth}
\begin{tabular}{lrr}
\toprule
Moment & Target & Model \\
\midrule
tfr & 1.700 & 1.898 \\
childless\_rate & 0.150 & 0.145 \\
mean\_age\_first\_birth & 26.000 & 33.535 \\
tfr\_gradient & 0.133 & 0.119 \\
own\_rate & 0.627 & 0.643 \\
own\_gradient & 0.170 & 0.139 \\
own\_family\_gap & 0.110 & 0.114 \\
prime\_childless\_renter\_median\_rooms & 4.000 & 6.365 \\
prime\_childless\_owner\_median\_rooms & 6.000 & 6.800 \\
housing\_increment\_0to1 & 0.664 & 0.441 \\
housing\_increment\_1to2 & 0.566 & 0.192 \\
young\_liquid\_wealth\_to\_income & 0.600 & 0.527 \\
center\_share\_nonparents & 0.494 & 0.405 \\
center\_share\_newparents & 0.416 & 0.382 \\
migration\_rate & 0.032 & 0.035 \\
old\_age\_own\_rate & 0.863 & 0.947 \\
old\_age\_parent\_childless\_gap & 0.070 & 0.062 \\
inv\_pop\_share\_C & 0.450 & 0.441 \\
inv\_rent\_ratio\_C\_over\_P & 1.140 & 1.182 \\
\bottomrule
\end{tabular}
\end{adjustbox}
\begin{minipage}{0.97\textwidth}\footnotesize \emph{Note:} These are the live benchmark moments from the current direct-geometry calibration record. Distributional tables below are not additional calibration targets.\end{minipage}
\end{table}


\begin{figure}[H]
\centering
\includegraphics[width=0.86\textwidth]{\detokenize{../output/model/distributional_discipline_current/report_figures/target_errors.pdf}}
\caption{Calibration target deviations. Positive values mean the model is above the target.}
\end{figure}

\section*{Income}

\begin{figure}[H]
\centering
\includegraphics[width=0.74\textwidth]{\detokenize{../output/model/distributional_discipline_current/report_figures/model_lifecycle_income_profile.pdf}}
\caption{Model lifecycle income profile. This is the deterministic age component that creates income dynamics in the model.}
\end{figure}

\begin{figure}[H]
\centering
\includegraphics[width=0.82\textwidth]{\detokenize{../output/model/distributional_discipline_current/report_figures/parenthood_by_income_bins.pdf}}
\caption{Parenthood by young-adult income bins. Data use PSID family income averaged over ages 25--30. The model uses ages 25--30 current income, which varies with age and location but not with idiosyncratic permanent income. Missing model points are empty rank bins.}
\end{figure}

\section*{Wealth And Tenure}

\begin{figure}[H]
\centering
\includegraphics[width=0.82\textwidth]{\detokenize{../output/model/distributional_discipline_current/report_figures/parenthood_by_wealth_bins.pdf}}
\caption{Parenthood by wealth bins. Data use young-adult PSID total wealth; model uses cross-sectional sale-net wealth at age 45.}
\end{figure}

\begin{figure}[H]
\centering
\includegraphics[width=0.82\textwidth]{\detokenize{../output/model/distributional_discipline_current/report_figures/childlessness_by_tenure_wealth.pdf}}
\caption{Childlessness by wealth and tenure. The data object follows young renters by whether they become owners by age 35; the model object is current tenure at age 45.}
\end{figure}

\section*{Housing Space}

\begin{figure}[H]
\centering
\includegraphics[width=0.86\textwidth]{\detokenize{../output/model/distributional_discipline_current/report_figures/rooms_mean_by_children.pdf}}
\caption{Mean rooms by tenure and children, ACS versus model.}
\end{figure}

\begin{figure}[H]
\centering
\includegraphics[width=0.86\textwidth]{\detokenize{../output/model/distributional_discipline_current/report_figures/rooms_share_ge7_by_children.pdf}}
\caption{Share of households with at least 7 rooms by tenure and children, ACS versus model.}
\end{figure}

\section*{Space And Mobility}

\begin{figure}[H]
\centering
\includegraphics[width=0.80\textwidth]{\detokenize{../output/model/distributional_discipline_current/report_figures/location_center_share.pdf}}
\caption{Center residence shares by parent status, MMS data versus model.}
\end{figure}

\begin{figure}[H]
\centering
\includegraphics[width=0.80\textwidth]{\detokenize{../output/model/distributional_discipline_current/report_figures/location_owner_rate.pdf}}
\caption{Owner rates by parent status, MMS data versus model.}
\end{figure}

\begin{figure}[H]
\centering
\includegraphics[width=0.80\textwidth]{\detokenize{../output/model/distributional_discipline_current/report_figures/location_rooms_by_parent.pdf}}
\caption{Mean rooms by parent status, MMS data versus model.}
\end{figure}

\begin{figure}[H]
\centering
\includegraphics[width=0.90\textwidth]{\detokenize{../output/model/distributional_discipline_current/report_figures/center_periphery_transitions.pdf}}
\caption{Center/periphery origin-destination shares. The model analog is the one-period location-choice probability averaged over the stationary distribution.}
\end{figure}

\begin{figure}[H]
\centering
\includegraphics[width=0.78\textwidth]{\detokenize{../output/model/distributional_discipline_current/report_figures/first_birth_mobility.pdf}}
\caption{First-birth moving facts. The current model has an analog for any location change, but no move-for-size or moved-to-ownership reason code.}
\end{figure}

\section*{Read}

\noindent The model has genuine lifecycle income dynamics, but not the cross-sectional income heterogeneity needed to match PSID income quintile fertility facts. The wealth comparison is closer, but still not exact because the data object is young-adult wealth predicting later fertility while the model object is a stationary cross-section. The room and spatial plots are the cleanest diagnostics: renter housing services are too high, the model does not generate enough room dispersion by tenure and children, and center/periphery transition rates are much too low. These are useful failures because they identify exactly which empirical margins the current model cannot yet claim to validate.

\end{document}
